% Springer Conference Paper Template
\documentclass{llncs}

% Packages
\usepackage{cite}
\usepackage{amsmath,amssymb,amsfonts}
\usepackage{algorithmic}
\usepackage{graphicx}
\usepackage{textcomp}
\usepackage{xcolor}
\usepackage{booktabs}
\usepackage{hyperref}
\setlength{\textfloatsep}{8pt plus 2pt minus 2pt}
\setlength{\floatsep}{8pt plus 2pt minus 2pt}
\setlength{\intextsep}{8pt plus 2pt minus 2pt}
\setlength{\abovecaptionskip}{4pt}
\setlength{\belowcaptionskip}{0pt}
\linespread{0.97}

\begin{document}

\title{A Hybrid Multi-Stage Plagiarism Detection Framework Combining Lexical and Semantic Similarity}

\author{Gurkirat Singh Bhangoo \and
Sparsh Tyagi \and
Ashmit Saini \and
Litesh Goyal \and
Ms. Ankita Thakur}

\institute{Apex Institute of Technology CSE, \\
Chandigarh University, Mohali, Punjab, India \\
\email{gurkirat.singh@email.com}, \email{sparshtyagi0612@gmail.com}, \email{sainiashmit.2000@gmail.com}, \email{goyallitesh195@gmail.com} (Orcid id-0009-0005-3795-0506), \email{ankita.e18488@cumail.in}}

\maketitle

\begin{abstract}
Plagiarism detection remains challenging because modern rewriting tools can preserve meaning while substantially altering surface wording. Traditional lexical approaches such as TF-IDF remain useful for exact or near-verbatim copying, but their performance drops when plagiarized text is paraphrased. Transformer-based semantic models such as Sentence-BERT (SBERT) help address this limitation, although they may also assign high similarity to documents that are topically related but independently written.

To address this trade-off, we present a hybrid multi-stage plagiarism detection framework that integrates exact matching, TF-IDF-based lexical similarity, SBERT-based semantic similarity, top-K sentence matching, candidate filtering, and document-level ensemble aggregation. The framework combines local sentence-level evidence with global document-level similarity signals, enabling detection of both direct copying and semantically preserved rewriting while reducing noisy matches.

We evaluate the framework on a controlled subset derived from the PAN plagiarism dataset using metadata-based positive pair construction and sampled negative pairs. On a held-out test split, the proposed hybrid model achieves precision of 0.9091, recall of 0.9524, F1-score of 0.9302, and accuracy of 0.9388, outperforming standalone TF-IDF and SBERT baselines. We also evaluate a logistic-regression classifier trained on the extracted hybrid features, report bootstrap confidence intervals, and conduct an ablation study. The results suggest that combining lexical and semantic cues at both local and global levels provides a strong and practically relevant basis for plagiarism detection on a controlled evaluation subset.

\keywords{plagiarism detection \and natural language processing \and TF-IDF \and Sentence-BERT \and semantic similarity \and hybrid classification systems \and information retrieval}
\end{abstract}

\section{Introduction}
Plagiarism detection remains an important problem in academic and digital environments, where large amounts of textual content are generated, shared, and repurposed online. Earlier systems mainly focused on exact copying, but modern plagiarism often involves paraphrasing, synonym substitution, and structural modification, which makes detection more difficult.

In our preliminary experiments, lexical methods such as TF-IDF worked well when copied passages preserved visible word overlap, but they were less effective for paraphrased content. At the same time, semantic models such as SBERT captured meaning more effectively, yet they sometimes marked unrelated sentences as similar because of shared context. These observations motivated us to develop a hybrid approach that uses both lexical and semantic similarity measures in a coordinated way.

The natural language processing community has responded with increasingly sophisticated semantic analysis tools. Transformer architectures, particularly BERT and its optimized variant Sentence-BERT, have demonstrated strong capabilities in understanding contextual meaning beyond surface-level lexical features. These models generate dense vector representations that capture semantic relationships, theoretically enabling detection even when every word has been changed. Yet this semantic sensitivity brings its own challenges. Documents discussing similar topics naturally exhibit semantic overlap, and distinguishing between legitimate topical similarity and actual content appropriation remains problematic.

We observe that lexical and semantic approaches each possess distinct strengths that could potentially address the other's weaknesses. Lexical methods provide high precision when matches occur, though their recall suffers with paraphrased content. Semantic methods cast a wider net, potentially catching paraphrased material, but risk flagging innocent topical similarities. Rather than choosing between these paradigms, our research explores how to effectively combine them within a unified detection pipeline.

This paper makes several contributions to plagiarism detection methodology. We introduce a multi-stage architecture that processes documents through exact matching, lexical scoring, semantic scoring, top-K sentence alignment, filtering, and document-level ensemble aggregation. We implement a top-K sentence matching strategy that reduces the cost of detailed pairwise comparison while preserving the strongest local evidence. We combine local sentence matches with global document-level TF-IDF and SBERT signals so that the final decision is not dominated by a single sentence pair. Finally, through empirical evaluation, we show that this hybrid approach produces stronger overall performance than either component method alone on the selected evaluation subset.

Our contribution is therefore best understood as a structured and reproducible hybrid framework rather than an entirely new similarity family. Recent literature already explores lexical-semantic combinations for plagiarism and text-similarity tasks; our emphasis is on a compact multi-stage design, careful threshold calibration, and an explicit analysis of which signals matter most on the chosen evaluation subset.

\section{Background and Related Research}

Classical plagiarism detection research established that paraphrasing is one of the main reasons lexical methods fail, motivating stronger semantic modeling and more fine-grained alignment analysis \cite{barronacevedo2013,foltynek2019}. TF-IDF-style models remain attractive because they are efficient, interpretable, and strong for direct copying, while transformer representations such as BERT and SBERT provide better coverage for semantically preserved rewriting \cite{devlin2019,reimers2019}. Recent comparative studies and surveys continue to report that no single family of methods is sufficient across all plagiarism types, especially when paraphrasing and AI-assisted rewriting are involved \cite{amirg2025,sajid2025,pudasaini2024}.

This has led to a growing body of hybrid work. Arabi and Akbari combine weighted lexical and semantic evidence for extrinsic plagiarism detection \cite{arabi2022}. A recent PAN 2025 participant system uses TF-IDF and Jaccard filtering before transformer classification \cite{luo2025}, while the PAN 2025 overview notes that many submissions relied on semantic alignment and that generalizability remains a challenge across datasets \cite{greinerpetter2025}. Concurrent lexical-semantic fusion studies outside plagiarism-specific benchmarks, such as TF-IDF + SBERT feature fusion for pairwise text classification, also report benefits from combining sparse and dense representations \cite{mehdi2025}. Our framework follows this broader direction, but focuses on a compact document-pair setting with top-K sentence matching, filtering, calibrated thresholds, ablation analysis, and explicit reporting of evaluation uncertainty.

\section{Proposed Methodology}

\subsection{Architecture Overview}

The pipeline begins with text normalization and sentence segmentation, followed by exact matching, document-level lexical similarity via TF-IDF, document-level semantic similarity via Sentence-BERT, and sentence-level semantic candidate generation. Rather than comparing all possible sentence pairs in full detail, we first rank sentence candidates by semantic similarity and retain only the top-K candidates for each source sentence. Candidate matches then undergo lexical-semantic filtering, and the surviving local evidence is combined with global document-level TF-IDF and SBERT scores to produce a final ensemble score. A tuned threshold converts this score into the final binary plagiarism decision.

This design uses cheap signals first, limits expensive pairwise comparison through top-K selection, and treats plagiarism as both a local alignment problem and a document-level similarity problem.

As shown in Fig.~\ref{fig:pipeline}, the proposed system follows a multi-stage pipeline integrating lexical and semantic similarity measures.
\begin{figure}[htbp]
\centerline{\includegraphics[width=0.38\textwidth]{pipeline.png}}
\caption{Proposed Hybrid Plagiarism Detection Pipeline}
\label{fig:pipeline}
\end{figure}

\subsection{Text Preprocessing}

Raw input documents are normalized by lowercasing, removing punctuation, and segmenting text into sentences.

\subsection{Exact Match Detection}

Before more expensive similarity calculations, we check whether any normalized sentences are exactly identical.

\subsection{Lexical Similarity via TF-IDF}

For non-identical text, lexical similarity is computed with TF-IDF cosine similarity.

\subsection{Semantic Similarity via Sentence-BERT}

To capture meaning beyond surface wording, we use Sentence-BERT with the all-MiniLM-L6-v2 model. Sentences are encoded into dense embeddings, and cosine similarity between embeddings measures semantic closeness.

\subsection{Top-K Sentence Matching}

Comparing every sentence pair is expensive, so we retain only the top-K semantic candidates for each source sentence.

\subsection{Adaptive Filtering Mechanism}

Weak sentence pairs are removed through lexical-semantic filtering.

\subsection{Score Aggregation Strategy}

Filtered matches are aggregated in two stages: a local hybrid score per sentence pair, followed by a document-level score combining local evidence with global TF-IDF and SBERT scores.

The aggregation function is:
$$\text{HybridScore} = w_e \cdot \mathbb{1}_{\text{exact}} + w_p \cdot S_{\text{TF-IDF}} + w_s \cdot S_{\text{SBERT}}$$

where $w_e$, $w_p$, and $w_s$ are tunable weights, $\mathbb{1}_{\text{exact}}$ is a binary indicator for exact matches, and the S terms represent normalized similarity scores.

Instead of reducing the document pair to a single maximum sentence score, we retain the strongest non-overlapping sentence matches and summarize them using peak local evidence, mean top-match evidence, matched sentence coverage, and exact-match ratio. We then combine this local signal with global lexical and semantic scores:

\begin{equation}
FinalScore = \alpha \cdot LocalSignal + \beta \cdot GlobalTFIDF + \gamma \cdot GlobalSBERT
\end{equation}

where $\alpha$, $\beta$, and $\gamma$ are tunable weights. This ensemble formulation makes the system more robust because a document-level plagiarism decision now depends on both local aligned evidence and whole-document similarity.

\subsection{Classification Decision}

The final step applies a threshold to the document-level hybrid score. Scores exceeding the threshold yield a positive plagiarism classification; scores below it indicate insufficient evidence of plagiarism. Rather than using one fixed threshold for every method, we tune thresholds independently for TF-IDF, SBERT, and the hybrid model on the training split.

Threshold choice reflects the usual precision-recall tradeoff.

\subsection{Learning-Based Decision Layer}

We additionally evaluate a logistic-regression classifier built on extracted hybrid features.

\section{Experimental Design}

We evaluate our proposed hybrid system through controlled experiments comparing its performance against baseline single-method approaches.

\subsection{Dataset Selection}

Our experiments utilize the PAN plagiarism detection dataset, a well-established benchmark in the research community. Instead of pairing documents by file order, we use the PAN metadata pair file to construct positive suspicious-source document pairs. Negative pairs are generated by sampling suspicious and source documents that are not linked in the PAN metadata. This yields a controlled dataset of 162 document pairs containing both plagiarized and non-plagiarized examples.

This construction is important because it ensures that positive labels reflect actual PAN pair relationships rather than positional assumptions in the file system. The resulting subset still represents only a portion of the full PAN corpus, but it provides a cleaner and more reproducible evaluation setup than the earlier experimental configuration.

\subsection{Evaluation Metrics}

We report precision, recall, F1-score, and accuracy on the held-out test set, together with bootstrap confidence intervals estimated from 500 resamples.

\subsection{Baseline Configurations}

We evaluate TF-IDF only, SBERT only, the proposed hybrid ensemble, and a learned classifier over extracted hybrid features. Each model receives an independently tuned threshold.

\subsection{Implementation Details}

We use a reproducible train-test protocol with a fixed random seed, giving 113 training pairs and 49 test pairs. The held-out test split contains 21 positive pairs and 28 negative pairs. Sentence-BERT uses all-MiniLM-L6-v2, sentence-level matching retains the top 3 semantic candidates per source sentence, and the learned classifier is implemented as logistic regression over extracted hybrid features.

\section{Results and Analysis}

Table \ref{tab:results} summarizes the main performance comparison.

\begin{table}[htbp]
\caption{Performance Comparison Across Detection Methods}
\begin{center}
\begin{tabular}{lcccc}
\toprule
\textbf{Method} & \textbf{Precision} & \textbf{Recall} & \textbf{F1} & \textbf{Accuracy} \\
\midrule
TF-IDF & 0.6897 & 0.9524 & 0.8000 & 0.7959 \\
SBERT & 0.7619 & 0.7619 & 0.7619 & 0.7959 \\
\textbf{Hybrid} & \textbf{0.9091} & \textbf{0.9524} & \textbf{0.9302} & \textbf{0.9388} \\
Learned Classifier & 0.9048 & 0.9048 & 0.9048 & 0.9184 \\
\bottomrule
\end{tabular}
\label{tab:results}
\end{center}
\end{table}
Table~\ref{tab:results} summarizes the comparative performance of all evaluated methods.

Fig.~\ref{fig:model} compares the evaluation metrics across the tested models.

\begin{figure}[htbp]
\centerline{\includegraphics[width=0.38\textwidth]{model_comparison.png}}
\caption{Comparison of Models based on F1 Score}
\label{fig:model}
\end{figure}

The hybrid model achieves the strongest overall performance, with precision of 0.9091, recall of 0.9524, F1-score of 0.9302, and accuracy of 0.9388. Its bootstrap confidence interval for F1 spans 0.8333 to 1.0000, which is encouraging but also reflects the modest test size. The learned classifier is competitive but does not surpass the tuned hybrid ensemble. TF-IDF remains strong on the present subset, while SBERT alone is weaker and less stable when used as the only signal.

\subsection{Ablation Analysis}

To better understand which parts of the proposed system contribute most to the final performance, we evaluate several ablation variants. These include a local-only model, a model without the local sentence-level signal, a model without global TF-IDF, and a model without global SBERT. The corresponding F1-scores are reported in Table~\ref{tab:ablation}.

\begin{table}[htbp]
\caption{Ablation Study on Hybrid Components}
\begin{center}
\begin{tabular}{lc}
\toprule
\textbf{Variant} & \textbf{F1 Score} \\
\midrule
Full Hybrid & \textbf{0.9302} \\
Minus Local Signal & 0.9302 \\
Minus Global TF-IDF & 0.8182 \\
Minus Global SBERT & 0.8235 \\
Local Only & 0.6000 \\
\bottomrule
\end{tabular}
\label{tab:ablation}
\end{center}
\end{table}

The ablation results show that the global semantic and lexical components contribute most of the discriminative power on the current split, while the local sentence-level signal adds comparatively little. This does not make the local stage useless in general; rather, it suggests that on this controlled subset the document-level signals already capture most of the separability. Since the held-out test set contains only 49 document pairs, the reported metrics should be interpreted as promising rather than definitive.

\section{Limitations and Future Directions}

Several limitations of our current work suggest directions for future research.

Our evaluation uses a controlled subset of 162 document pairs and only 49 held-out test pairs, so the results should be interpreted cautiously. In addition, the negative pairs are sampled rather than retrieved from a realistic large-scale candidate pool, which may make the task easier than a full operational setting. The study also does not reproduce the full PAN 2025 retrieval-and-alignment setup, so the reported numbers should not be interpreted as directly comparable with shared-task leaderboards. Future work should therefore evaluate the method on larger benchmarks, harder negatives, multilingual plagiarism, and more deployment-oriented retrieval scenarios.

\section{Conclusion}

This work demonstrates that carefully designed hybrid approaches can substantially outperform single-method baselines for plagiarism detection. By combining exact matching, document-level lexical similarity, document-level semantic similarity, and filtered top-K sentence evidence, the proposed system achieves precision of 0.9091, recall of 0.9524, F1-score of 0.9302, and accuracy of 0.9388 on a held-out evaluation split derived from PAN metadata.

The results support the use of a hybrid lexical-semantic framework for plagiarism detection on controlled document-pair evaluations. The proposed system improves over standalone baselines, remains competitive even when compared with a learned classifier, and provides a practical basis for future plagiarism-screening systems. At the same time, the controlled evaluation setting, sampled negatives, and limited test size indicate that broader validation is still needed before stronger general claims are made.

\begin{thebibliography}{99}

\bibitem{potthast2010}
M. Potthast, B. Stein, A. Barrón-Cedeño, and P. Rosso, ``An evaluation framework for plagiarism detection,'' in \textit{Proceedings of the 23rd International Conference on Computational Linguistics}, Beijing, China, 2010, pp. 997--1005.

\bibitem{reimers2019}
N. Reimers and I. Gurevych, ``Sentence-BERT: Sentence embeddings using siamese BERT-networks,'' in \textit{Proceedings of the 2019 Conference on Empirical Methods in Natural Language Processing}, Hong Kong, China, 2019, pp. 3982--3992.

\bibitem{salton1988}
G. Salton and C. Buckley, ``Term-weighting approaches in automatic text retrieval,'' \textit{Information Processing \& Management}, vol. 24, no. 5, pp. 513--523, 1988.

\bibitem{devlin2019}
J. Devlin, M.-W. Chang, K. Lee, and K. Toutanova, ``BERT: Pre-training of deep bidirectional transformers for language understanding,'' in \textit{Proceedings of the 2019 Conference of the North American Chapter of the Association for Computational Linguistics}, Minneapolis, MN, USA, 2019, pp. 4171--4186.

\bibitem{mikolov2013}
T. Mikolov, I. Sutskever, K. Chen, G. Corrado, and J. Dean, ``Distributed representations of words and phrases and their compositionality,'' in \textit{Advances in Neural Information Processing Systems 26}, Lake Tahoe, NV, USA, 2013, pp. 3111--3119.

\bibitem{pennington2014}
J. Pennington, R. Socher, and C. Manning, ``GloVe: Global vectors for word representation,'' in \textit{Proceedings of the 2014 Conference on Empirical Methods in Natural Language Processing}, Doha, Qatar, 2014, pp. 1532--1543.

\bibitem{vaswani2017}
A. Vaswani, N. Shazeer, N. Parmar, J. Uszkoreit, L. Jones, A. N. Gomez, Ł. Kaiser, and I. Polosukhin, ``Attention is all you need,'' in \textit{Advances in Neural Information Processing Systems 30}, Long Beach, CA, USA, 2017, pp. 5998--6008.

\bibitem{barronacevedo2013}
A. Barrón-Cedeño, M. Vila, M. A. Martí, and P. Rosso, ``Plagiarism meets paraphrasing: Insights for the next generation in automatic plagiarism detection,'' \textit{Computational Linguistics}, vol. 39, no. 4, pp. 917--947, 2013.

\bibitem{foltynek2019}
T. Foltýnek, N. Meuschke, and B. Gipp, ``Academic plagiarism detection: A systematic literature review,'' \textit{ACM Computing Surveys}, vol. 52, no. 6, pp. 112:1--112:42, 2019.

\bibitem{arabi2022}
H. Arabi and M. Akbari, ``Improving plagiarism detection in text document using hybrid weighted similarity,'' \textit{Expert Systems with Applications}, vol. 207, p. 118034, 2022.

\bibitem{pudasaini2024}
S. Pudasaini, L. Miralles-Pechu\'an, D. Lillis, and M. Llorens Salvador, ``Survey on plagiarism detection in large language models: The impact of ChatGPT and Gemini on academic integrity,'' arXiv:2407.13105, 2024.

\bibitem{greinerpetter2025}
A. Greiner-Petter, M. Fr\"obe, J. P. Wahle, T. Ruas, B. Gipp, A. Aizawa, and M. Potthast, ``Overview of the plagiarism detection task at PAN 2025,'' in \textit{CLEF 2025 Working Notes}, CEUR Workshop Proceedings, vol. 4038, Madrid, Spain, 2025, pp. 3575--3585.

\bibitem{luo2025}
J. Luo, M. Huang, B. Liu, and Z. Han, ``Team jr at generative plagiarism detection 2025: A two-stage approach from TF-IDF/Jaccard filtering to transformer classification,'' in \textit{CLEF 2025 Working Notes}, CEUR Workshop Proceedings, vol. 4038, Madrid, Spain, 2025.

\bibitem{mehdi2025}
M. Mehdi, S. Mushtaq, and G. R. Butt, ``A hybrid TF-IDF and SBERT approach for enhanced text classification performance,'' \textit{Preprints.org}, 2025. doi:10.20944/preprints202510.2427.v1.

\bibitem{amirg2025}
A. Amirzhanov, C. Turan, and A. Makhmutova, ``Plagiarism types and detection methods: A systematic survey of algorithms in text analysis,'' \textit{Frontiers in Computer Science}, vol. 7, p. 1504725, 2025.

\bibitem{sajid2025}
M. Sajid, M. Sanaullah, M. Fuzail, T. S. Malik, and S. M. Shuhidan, ``Comparative analysis of text-based plagiarism detection techniques,'' \textit{PLOS ONE}, vol. 20, no. 4, e0319551, 2025.

\end{thebibliography}

\end{document}
